\documentclass[11pt,a4paper]{article}
\usepackage[T1]{fontenc}
\usepackage[utf8]{inputenc}
\usepackage[turkish]{babel}
\usepackage{microtype}
\usepackage{hyperref}

\title{\textbf{Hidrasyon Hatası -- Ne Anlama Geliyor?}}
\author{}
\date{}

\begin{document}
\maketitle

\section*{Özet}
SSR (\textit{Server-Side Rendering}) ile gönderilen HTML, tarayıcıda React (veya Next.js) tarafından \emph{hidrate} edilir; yani DOM, React bileşen ağacına yeniden bağlanır.

\section*{Hata ne zaman oluşur?}
\textbf{Hydration failed} (hidrasyon başarısız) hatası, sunucunun ürettiği HTML ile istemcinin (tarayıcının) aynı bileşeni \textbf{ilk render} ettiğinde ürettiği HTML'in \textbf{uyuşmaması} durumunda ortaya çıkar.

\end{document}
